\section{Premier Modèle}
\label{sec:PremierModel}

Je me suis d'abord familiarisé avec Java et la PPC en réalisant un
premier modèle pour répondre à un premier problème :
\[\text{A partir d'un accord, trouver toutes les positions jouables sur un 
manche de guitare.}\]

Pour ce problème, nous modélisons dans un premier temps les objets
étudiés (guitare, Guitariste, partition) puis nous décrivons des contraintes
qui assurent une résolution du problème.

\subsection{Modélisation des objets}
\label{subsec:ModelGuitare}

On implémente une guitare par :
\begin{itemize}
  \item \nbCordes : nombre de cordes (entier)
  \item \nbFrettes : nombre de frettes (entier)
  \item \scordature : scordature de la guitare (tableau d'entiers)
\end{itemize}

On implémente un ou une guitariste par :
\begin{itemize}
  \item \tailleMain : la taille de la main en le nombre de frêtes accessibles
    depuis le haut du manche (entier)
	\item \nbDoigtsUtil : le nombre de doigts utilisés pour jouer (entier)
	\item \isBarre : $-1$ s'il est impossible de réaliser un barré, le
  nombre de frettes accessibles quand on fait un barré sinon (entier)
\end{itemize}

\subsection{La distance sur le manche}
\label{subse:distanceManche}

Sur le manche de la guitare, les frettes ne sont pas disposées de manière
régulière car le demi ton n'est pas défini de manière linéaire.
Par exemple, pour augmenter d'une octave, on divise la longueur de la
corde par deux. On peut alors trouver une relation entre le nombre de frettes
accessibles depuis une position et \tailleMain.

\begin{equation}
  \label{eq:taille_frette}
  \dLegale(position) = \tailleMain\times 2^{\frac{position}{12}}
\end{equation}

\subsection{Modèle PPC}
\label{subsec:ModelePPC}

Nous définisons ici le premier modèle PPC réalisé.

\begin{defi}(Indéterminées)
  \label{def:IndetPremier}
  Les indéterminées du modèle sont les frettes pressées sur chaque corde. Cela
  permet donc de définir une position
  $\indetPremier = \accolades{\indetPremier_i}_{i\in I}$ où $I$ est une
  numérotation des cordes.
  Il y a donc \nbCordes \ indéterminées à valeur dans $\internat{-1, \nbFrettes}$.
  $-1$ représente une corde non jouée, sinon c'est une corde jouée et pressée
  sur la frette associée. $0$ représente la corde à vide.
\end{defi}

\begin{defi}(Contraintes)
  \label{def:ContraintesPremier}
  Pour ce modèle, nous avons un ensemble de contraintes. Premièrement, il faut
  que les notes jouées soient uniquement des notes de l'accord. Cela nous
  amène à la première contrainte :
  \begin{equation}
    \label{eq:played_in_chord}
    \forall i\in \internat{ 1, \nbCordes }, \exists j\geq 1,
    (\scordature[i]+\indetPremier_i)-(\accord[0]+\accord[j])
    = 0[\mod 12]
  \end{equation}

  Il faut aussi que l'ensemble des notes de l'accord soient jouées, cela nous
  amène à la deuxième contrainte :
  \begin{equation}
    \label{eq:chord_play}
      \forall i\in \internat{ 1,\lenAccord }, \exists j\geq 1,
      (\scordature[j]+\indetPremier_j)-(\accord[0]+\accord[i])
    = 0[\mod 12]
  \end{equation}
  
  On rappelle que l'octave n'est pas importante et qu'il y a $12$ demis tons
  dans une octave d'où le modulo $12$.\\
  Nous avons pour le moment des contraintes permettant d'assurer que ce que
  l'on joue correspond à la suite d'accords.\\
  Il faut de plus modéliser la contrainte de distance pour avoir une première
  notion de jouabilité de la position.
  Il faut que la distance entre les doigts ne dépasse pas une certainte borne.
\end{defi}

Après implémentation, le solver choco trouve toutes les solutions du modèle
demandé, donc toutes les positions qui respectent les contraintes pour un 
accord donné.

Nous avons donc créé un modèle de
contraintes permettant, à partir du chiffrage d'un accord, de trouver
l'ensemble des positions physiquement jouables de cet accord.
Cependant la notion de physiquement jouable reste assez limitée.
Et on s'est pour l'instant restraint à un seul accord.
\newpage
